% Options for packages loaded elsewhere
\PassOptionsToPackage{unicode}{hyperref}
\PassOptionsToPackage{hyphens}{url}
\documentclass[
  british,
]{article}
\usepackage{xcolor}
\usepackage{amsmath,amssymb}
\setcounter{secnumdepth}{5}
\usepackage{iftex}
\ifPDFTeX
  \usepackage[T1]{fontenc}
  \usepackage[utf8]{inputenc}
  \usepackage{textcomp} % provide euro and other symbols
\else % if luatex or xetex
  \usepackage{unicode-math} % this also loads fontspec
  \defaultfontfeatures{Scale=MatchLowercase}
  \defaultfontfeatures[\rmfamily]{Ligatures=TeX,Scale=1}
\fi
\usepackage{lmodern}
\ifPDFTeX\else
  % xetex/luatex font selection
\fi
% Use upquote if available, for straight quotes in verbatim environments
\IfFileExists{upquote.sty}{\usepackage{upquote}}{}
\IfFileExists{microtype.sty}{% use microtype if available
  \usepackage[]{microtype}
  \UseMicrotypeSet[protrusion]{basicmath} % disable protrusion for tt fonts
}{}
\makeatletter
\@ifundefined{KOMAClassName}{% if non-KOMA class
  \IfFileExists{parskip.sty}{%
    \usepackage{parskip}
  }{% else
    \setlength{\parindent}{0pt}
    \setlength{\parskip}{6pt plus 2pt minus 1pt}}
}{% if KOMA class
  \KOMAoptions{parskip=half}}
\makeatother
\usepackage{color}
\usepackage{fancyvrb}
\newcommand{\VerbBar}{|}
\newcommand{\VERB}{\Verb[commandchars=\\\{\}]}
\DefineVerbatimEnvironment{Highlighting}{Verbatim}{commandchars=\\\{\}}
% Add ',fontsize=\small' for more characters per line
\newenvironment{Shaded}{}{}
\newcommand{\AlertTok}[1]{\textcolor[rgb]{1.00,0.00,0.00}{\textbf{#1}}}
\newcommand{\AnnotationTok}[1]{\textcolor[rgb]{0.38,0.63,0.69}{\textbf{\textit{#1}}}}
\newcommand{\AttributeTok}[1]{\textcolor[rgb]{0.49,0.56,0.16}{#1}}
\newcommand{\BaseNTok}[1]{\textcolor[rgb]{0.25,0.63,0.44}{#1}}
\newcommand{\BuiltInTok}[1]{\textcolor[rgb]{0.00,0.50,0.00}{#1}}
\newcommand{\CharTok}[1]{\textcolor[rgb]{0.25,0.44,0.63}{#1}}
\newcommand{\CommentTok}[1]{\textcolor[rgb]{0.38,0.63,0.69}{\textit{#1}}}
\newcommand{\CommentVarTok}[1]{\textcolor[rgb]{0.38,0.63,0.69}{\textbf{\textit{#1}}}}
\newcommand{\ConstantTok}[1]{\textcolor[rgb]{0.53,0.00,0.00}{#1}}
\newcommand{\ControlFlowTok}[1]{\textcolor[rgb]{0.00,0.44,0.13}{\textbf{#1}}}
\newcommand{\DataTypeTok}[1]{\textcolor[rgb]{0.56,0.13,0.00}{#1}}
\newcommand{\DecValTok}[1]{\textcolor[rgb]{0.25,0.63,0.44}{#1}}
\newcommand{\DocumentationTok}[1]{\textcolor[rgb]{0.73,0.13,0.13}{\textit{#1}}}
\newcommand{\ErrorTok}[1]{\textcolor[rgb]{1.00,0.00,0.00}{\textbf{#1}}}
\newcommand{\ExtensionTok}[1]{#1}
\newcommand{\FloatTok}[1]{\textcolor[rgb]{0.25,0.63,0.44}{#1}}
\newcommand{\FunctionTok}[1]{\textcolor[rgb]{0.02,0.16,0.49}{#1}}
\newcommand{\ImportTok}[1]{\textcolor[rgb]{0.00,0.50,0.00}{\textbf{#1}}}
\newcommand{\InformationTok}[1]{\textcolor[rgb]{0.38,0.63,0.69}{\textbf{\textit{#1}}}}
\newcommand{\KeywordTok}[1]{\textcolor[rgb]{0.00,0.44,0.13}{\textbf{#1}}}
\newcommand{\NormalTok}[1]{#1}
\newcommand{\OperatorTok}[1]{\textcolor[rgb]{0.40,0.40,0.40}{#1}}
\newcommand{\OtherTok}[1]{\textcolor[rgb]{0.00,0.44,0.13}{#1}}
\newcommand{\PreprocessorTok}[1]{\textcolor[rgb]{0.74,0.48,0.00}{#1}}
\newcommand{\RegionMarkerTok}[1]{#1}
\newcommand{\SpecialCharTok}[1]{\textcolor[rgb]{0.25,0.44,0.63}{#1}}
\newcommand{\SpecialStringTok}[1]{\textcolor[rgb]{0.73,0.40,0.53}{#1}}
\newcommand{\StringTok}[1]{\textcolor[rgb]{0.25,0.44,0.63}{#1}}
\newcommand{\VariableTok}[1]{\textcolor[rgb]{0.10,0.09,0.49}{#1}}
\newcommand{\VerbatimStringTok}[1]{\textcolor[rgb]{0.25,0.44,0.63}{#1}}
\newcommand{\WarningTok}[1]{\textcolor[rgb]{0.38,0.63,0.69}{\textbf{\textit{#1}}}}
\usepackage{longtable,booktabs,array}
\usepackage{caption}
\captionsetup[table]{skip=6pt}
\newcounter{none} % for unnumbered tables
\usepackage{calc} % for calculating minipage widths
% Correct order of tables after \paragraph or \subparagraph
\usepackage{etoolbox}
\makeatletter
\patchcmd\longtable{\par}{\if@noskipsec\mbox{}\fi\par}{}{}
\makeatother
\usepackage{graphicx}
\makeatletter
\newsavebox\pandoc@box
\newcommand*\pandocbounded[1]{% scales image to fit in text height/width
  \sbox\pandoc@box{#1}%
  \Gscale@div\@tempa{\textheight}{\dimexpr\ht\pandoc@box+\dp\pandoc@box\relax}%
  \Gscale@div\@tempb{\linewidth}{\wd\pandoc@box}%
  \ifdim\@tempb\p@<\@tempa\p@\let\@tempa\@tempb\fi% select the smaller of both
  \ifdim\@tempa\p@<\p@\scalebox{\@tempa}{\usebox\pandoc@box}%
  \else\usebox{\pandoc@box}%
  \fi%
}
% Set default figure placement to htbp
\def\fps@figure{htbp}
\makeatother
% definitions for citeproc citations
\NewDocumentCommand\citeproctext{}{}
\NewDocumentCommand\citeproc{mm}{%
  \begingroup\def\citeproctext{#2}\cite{#1}\endgroup}
\makeatletter
 % allow citations to break across lines
 \let\@cite@ofmt\@firstofone
 % avoid brackets around text for \cite:
 \def\@biblabel#1{}
 \def\@cite#1#2{{#1\if@tempswa , #2\fi}}
\makeatother
\newlength{\cslhangindent}
\setlength{\cslhangindent}{1.5em}
\newlength{\csllabelwidth}
\setlength{\csllabelwidth}{3em}
\newenvironment{CSLReferences}[2] % #1 hanging-indent, #2 entry-spacing
 {\begin{list}{}{%
  \setlength{\itemindent}{0pt}
  \setlength{\leftmargin}{0pt}
  \setlength{\parsep}{0pt}
  % turn on hanging indent if param 1 is 1
  \ifodd #1
   \setlength{\leftmargin}{\cslhangindent}
   \setlength{\itemindent}{-1\cslhangindent}
  \fi
  % set entry spacing
  \setlength{\itemsep}{#2\baselineskip}}}
 {\end{list}}
\usepackage{calc}
\newcommand{\CSLBlock}[1]{\hfill\break\parbox[t]{\linewidth}{\strut\ignorespaces#1\strut}}
\newcommand{\CSLLeftMargin}[1]{\parbox[t]{\csllabelwidth}{\strut#1\strut}}
\newcommand{\CSLRightInline}[1]{\parbox[t]{\linewidth - \csllabelwidth}{\strut#1\strut}}
\newcommand{\CSLIndent}[1]{\hspace{\cslhangindent}#1}
\ifLuaTeX
\usepackage[bidi=basic,shorthands=off]{babel}
\else
\usepackage[bidi=default,shorthands=off]{babel}
\fi
\ifLuaTeX
  \usepackage{selnolig} % disable illegal ligatures
\fi
\setlength{\emergencystretch}{3em} % prevent overfull lines
\providecommand{\tightlist}{%
  \setlength{\itemsep}{0pt}\setlength{\parskip}{0pt}}
\usepackage{bookmark}
\IfFileExists{xurl.sty}{\usepackage{xurl}}{} % add URL line breaks if available
\urlstyle{same}
% fallback for those not using the hyperref driver hyperxmp:
\makeatletter
\@ifundefined{xmpquote}{\newcommand{\xmpquote}[1]{#1}}{}
\makeatother
\hypersetup{
  pdftitle={Coarse-to-Fine Text Generation on a Laptop: A Native MLX Study of Next-Concept Modeling},
  pdfauthor={Jan Music},
  pdflang={en-GB},
  pdfkeywords={\xmpquote{coarse-to-fine generation}, \xmpquote{residual
vector quantisation}, \xmpquote{parallel decoding}, \xmpquote{Apple
Silicon}, \xmpquote{MLX}},
  hidelinks,
  pdfcreator={LaTeX via pandoc}}

\title{Coarse-to-Fine Text Generation on a Laptop: A Native MLX Study of
Next-Concept Modeling}
\author{Jan Music}
\date{v0.11 · 2026-10-27}

\begin{document}
\maketitle
\begin{abstract}
Next-Concept Modeling (NCM) generates a sequence through successive
levels of discrete latent codes, predicting each level in parallel. We
independently implement the method in MLX and study synthetic language
and Sudoku on Apple Silicon and NVIDIA GPUs. On a finite synthetic
grammar, a 5M-parameter generator achieves lower small-batch latency
than a similarly sized autoregressive generator, with lower sentence
validity (89.8\% versus 98.6\%). Two Sudoku training seeds use the same
6.44M-parameter generator recipe and 16.384M presentations each.
Retained greedy evaluation aggregates report easy/medium/hard solve
rates of 98.10/83.40/40.30\% for seed 0 and 97.25/79.60/33.25\% for seed
1, with 2,000 boards per difficulty. The hard-set difference is 7.05
percentage points; two seeds do not establish population training
variability or its cause. Original NCM reports 99.50/96.50/86.10\% with
a 25.46M generator, 51.2M presentations and a different task setup. Our
generator is 3.95 times smaller, excluding its additional codec; the
comparison establishes neither task parity nor matched-cost efficiency.
E20 completes its schedule without a recorded reset or truncation, but
missing per-board predictions and remote provenance manifests limit
independent audit. We characterize this phase as a partial
implementation and small-model success, retain the substantial
hard-Sudoku deficit, and define controlled text evaluation as the next
research direction rather than a demonstrated capability.
\end{abstract}

\begin{quote}
\textbf{Status.} Working paper, version 0.11. Every number here is
copied from a run in the code repository (\texttt{runs/}, summarised in
\texttt{paper/data/results.json}) or quoted from the cited paper and
marked as quoted. Experiments still running are marked \emph{in
progress} and never estimated. Appendix C lists the changes between
versions.
\end{quote}

\section{Reading guide}\label{reading-guide}

This paper has two audiences. \textbf{If you are new to language
models}, read Section 2 (the idea in plain language) and Section 3
(glossary), then the boxed summary at the start of each later section.
\textbf{If you build models}, Sections 5--7 give the method with every
formula explained, the Apple Silicon implementation and the experiments.
Section 8 tells the story of the project run by run: what each
experiment (E0, E1, \ldots) asked, what changed and what we learned.

\section{The idea in plain language}\label{the-idea-in-plain-language}

\subsection{How a chat model writes}\label{how-a-chat-model-writes}

A chat model such as ChatGPT or Qwen writes one \emph{token} (a word or
a piece of a word) at a time. To write the 100th token it runs the whole
model once, reads what it has written so far and picks the next token. A
100-token answer therefore needs 100 full runs of the model, one after
the other. At small batch sizes, autoregressive decoding often spends
much of its time moving model weights from memory. The balance depends
on the model, hardware, cache and implementation; these experiments
measure latency, not memory-bandwidth utilisation.

\subsection{How next-concept modeling
writes}\label{how-next-concept-modeling-writes}

NCM (\citeproc{ref-ncm2026}{Kumbong et al. 2026}) writes differently. It
first produces a single, very rough ``concept'' for the whole text, then
refines it in a few rounds. Each round adds roughly twice as much
detail, and \textbf{every piece of detail in a round is produced at the
same time}. Figure 1 shows a real example from this project: a
32-character sentence written in six rounds (levels). After level 0 the
rough shape is there (``eue likes one oed egg.''); by level 2 the words
are settled; later levels only change single letters.

\begin{figure}
\centering
\includegraphics[width=1\linewidth,height=\textheight,keepaspectratio,alt={A
real generation from experiment E1. Left: the codes of each level; all
codes in one row come from one model run. Right: the text obtained by
decoding everything up to that level. Six model runs write the sentence;
a word-by-word model needs 32. The sample also shows a failure: level 3
turns ``red'' into ``ret'', and nothing later repairs
it.}]{figures/overview.pdf}
\caption{A real generation from experiment E1. Left: the codes of each
level; all codes in one row come from one model run. Right: the text
obtained by decoding everything up to that level. Six model runs write
the sentence; a word-by-word model needs 32. The sample also shows a
failure: level 3 turns ``red'' into ``ret'', and nothing later repairs
it.}
\end{figure}

The number of rounds is fixed by the model configuration (six in
TinyWorld, twelve in the original paper's Sudoku model), so NCM needs
far fewer model runs than a word-by-word model. That is where its speed
comes from.

\subsection{Two parts: a codec and a
generator}\label{two-parts-a-codec-and-a-generator}

NCM has two learned parts. An analogy from image files helps.

\begin{itemize}
\tightlist
\item
  The \textbf{codec} is like an image format with progressive loading: a
  JPEG that first shows a blurry picture, then sharpens. It learns to
  turn a sentence into a stack of codes, from one coarse code to one
  code per letter, and to turn the codes back into the sentence. It is
  trained first, on its own, with reconstruction losses. Full-hierarchy
  reconstruction is measured separately from generation quality.
\item
  The \textbf{generator} is the creative part. It learns to
  \emph{invent} new code stacks one level at a time: level 0, then level
  1 given level 0, and so on. The codec then turns the finished stack
  into text.
\end{itemize}

\subsection{Why it matters, and the
catch}\label{why-it-matters-and-the-catch}

Fewer, wider model runs suit a laptop GPU: one run reads the weights
once and does a lot of arithmetic with them. The catch is that all codes
of a level are chosen at the same moment, so neighbouring pieces cannot
see each other's choice. If two pieces of a level disagree, the result
can be wrong in a way a word-by-word model would avoid. Much of this
paper measures when that matters.

\subsection{What we did}\label{what-we-did}

We rebuilt NCM from scratch for Apple's MLX framework
(\citeproc{ref-mlx2023}{Hannun et al. 2023}), tested it on a small
invented language and on Sudoku puzzles (the original paper's showcase),
and compared it with a character-level AR model of similar generator
size and with a large chat model. The speed advantage holds on a laptop.
On Sudoku, two training seeds solve 97.25--98.10\% of easy boards but
only 33.25--40.30\% of hard boards. Original NCM reports 99.5\% and
86.1\%, respectively, under a different setup. A historical comparison
matches generator presentations against our AR baseline; Section 8
records the experiments. NCM is a generation method used here for two
tasks, not a Sudoku-specific expansion of the term or an
instruction-following language model.

\section{Glossary}\label{glossary}

\begin{longtable}[]{@{}
  >{\raggedright\arraybackslash}p{(\linewidth - 2\tabcolsep) * \real{0.5000}}
  >{\raggedright\arraybackslash}p{(\linewidth - 2\tabcolsep) * \real{0.5000}}@{}}
\caption{Terms used in this paper.}\tabularnewline
\toprule\noalign{}
\begin{minipage}[b]{\linewidth}\raggedright
Term
\end{minipage} & \begin{minipage}[b]{\linewidth}\raggedright
Meaning
\end{minipage} \\
\midrule\noalign{}
\endfirsthead
\toprule\noalign{}
\begin{minipage}[b]{\linewidth}\raggedright
Term
\end{minipage} & \begin{minipage}[b]{\linewidth}\raggedright
Meaning
\end{minipage} \\
\midrule\noalign{}
\endhead
\bottomrule\noalign{}
\endlastfoot
Token & The unit a model reads and writes. Here: one character
(TinyWorld) or one Sudoku cell. \\
Autoregressive (AR) model & A model that writes one token at a time,
each conditioned on all previous ones. Chat models are AR models. \\
Forward pass, NFE & One run of a model on its input. NFE (number of
function evaluations) counts forward passes per generated sample. \\
Parameters & The learned numbers inside a model; ``6.4M parameters''
means 6.4 million. \\
Embedding, latent vector & A list of numbers (128--512 here) that
represents a token or a position inside the model. \\
Vector quantisation (VQ) & Replacing a vector by the nearest entry of a
fixed list, the \emph{codebook}, so it can be stored as one integer, the
\emph{code}. \\
Codebook, code & The list of allowed vectors at one level (16 or 128
entries here) and the index of one entry. \\
Level & One round of refinement. Level \(k\) has \(l_k\) codes; level 0
is the coarsest. \\
Residual & What is left over: the difference between the true vector and
the approximation built so far. \\
Codec (tokenizer, VQ-VAE) & Stage 1 of NCM: encoder, quantiser and
decoder that turn text into codes and back. \\
Generator & Stage 2 of NCM: the transformer that predicts the codes of
each level. \\
Transformer, attention, mask & The standard neural network for
sequences. \emph{Attention} lets each position read information from
other positions; a \emph{mask} says which positions may be read. \\
Teacher forcing & Training on the \emph{true} previous outputs instead
of the model's own. Fast, but the model never practises recovering from
its own mistakes. \\
Conditioning corruption & Replacing some true inputs with wrong ones
during training so the model learns to cope with its own mistakes. \\
Cross-entropy (CE), nats & The training loss: how surprised the model is
by the right answer, in natural-log units. Lower is better;
\(\ln 16 \approx 2.77\) nats is a uniform guess among 16 codes. \\
Perplexity & \(e^{\mathrm{CE}}\), the effective number of choices. Code
perplexity 1 means a level always uses the same code and carries no
information. \\
Temperature, greedy & Sampling randomness. Temperature 0 (greedy) always
takes the most likely choice. \\
Batch, latency, p50 & A batch is several samples generated together;
latency is the time for one batch; p50 is the median over repeated
runs. \\
KV cache & A store of past attention values that speeds up AR
generation. \\
EMA & Exponential moving average: a running average that slowly forgets
old values. Used to update codebooks. \\
Straight-through estimator & A trick to train through a step that has no
gradient (like picking the nearest code) by treating it, for the
gradient only, as the identity. \\
MLX, Metal & Apple's machine-learning array framework, and the GPU
interface it runs on. \\
fp32, bf16 & 32-bit and 16-bit (bfloat16) floating-point formats. \\
E0, E1, \ldots{} & Experiment numbers, in the order the experiments were
run (Section 8). \\
\end{longtable}

\section{Background}\label{background}

\textbf{Residual vector quantisation.} A vector quantiser maps a
continuous vector to the nearest of \(V\) codebook entries
(\citeproc{ref-vqvae2017}{Oord et al. 2017}). Residual quantisation
applies several quantisers in sequence, each to what the previous ones
left over (\citeproc{ref-rqvae2022}{Lee et al. 2022};
\citeproc{ref-soundstream2021}{Zeghidour et al. 2021}); it underlies
modern neural audio codecs.

\textbf{Next-scale prediction.} VAR (\citeproc{ref-var2024}{Tian et al.
2024}) combines residual quantisation with resolution: the first
quantiser sees a \(1\times1\) summary of an image, later ones
\(2\times2\), \(3\times3\), up to full resolution, and a transformer
generates one scale per step. NCM applies the idea to 1-D token
sequences (\citeproc{ref-ncm2026}{Kumbong et al. 2026}).

\textbf{Diffusion language models} refine positions of a sequence over
\(T\) denoising steps (\citeproc{ref-mdlm2024}{Sahoo et al. 2024}).
Full-sequence samplers process length \(L\) at each step; NCM's early
levels process fewer code positions, although prefix and decoder work
also contribute. The NCM paper reports quality at or above strong
diffusion language models with 10--31× faster generation under its H100
benchmark protocol (\citeproc{ref-ncm2026}{Kumbong et al. 2026}).

\section{Method}\label{method}

\begin{quote}
\textbf{In short.} The codec learns to describe a sequence with a stack
of codes, coarse to fine. The generator learns to predict each level's
codes from the levels above it. Generation runs the generator once per
level, then decodes.
\end{quote}

We follow the NCM design (\citeproc{ref-ncm2026}{Kumbong et al. 2026})
and its released Sudoku configuration; Table 2 lists where we differ. A
running toy example uses a sequence of \(L = 8\) tokens, levels
\(l = (1, 2, 4, 8)\) and vectors of width \(C = 2\).

\subsection{Codec: from tokens to a stack of
codes}\label{codec-from-tokens-to-a-stack-of-codes}

\textbf{Encoder.} Each token \(x_i\) and its position \(i\) are looked
up in two tables, added and normalised:

\[z_i = \mathrm{Norm}\big(W_{\mathrm{tok}}[x_i] + W_{\mathrm{pos}}[i]\big), \qquad z \in \mathbb{R}^{L \times C}.\]

\emph{Reading it:} \(W_{\mathrm{tok}}\) has one learned row per possible
token and \(W_{\mathrm{pos}}\) one per position; \(z\) is the whole
sequence as \(L\) vectors of \(C\) numbers. ``Norm'' rescales each
vector (LayerNorm), or groups of channels across the sequence
(GroupNorm, as upstream), without a learned scale, because the quantiser
compares squared distances and a learned scale could grow without bound.
In the toy example \(z\) is an \(8 \times 2\) table of numbers.

\textbf{Quantisation, level by level.} The quantiser keeps a
\emph{residual} \(r\), initially all of \(z\), and for each level
\(k = 0, \dots, K-1\) does four things:

\[\underbrace{x_k = D_k\, r_k}_{\text{1. shrink}} ,\qquad
\underbrace{c_k = \arg\min_j \lVert x_k - E_k[j] \rVert^2}_{\text{2. pick nearest codes}},\]
\[\underbrace{q_k = \phi_k\big(U_k\, E_k[c_k]\big)}_{\text{3. expand and refine}},\qquad
\underbrace{r_{k+1} = r_k - \mathrm{sg}(q_k)}_{\text{4. subtract}} .\]

\emph{Reading it:}

\begin{enumerate}
\def\labelenumi{\arabic{enumi}.}
\tightlist
\item
  \textbf{Shrink.} \(D_k\) averages the residual down to \(l_k\)
  positions (area pooling). Toy example: for level 1 (\(l_1 = 2\)),
  positions 1--4 are averaged into one vector and 5--8 into another.
\item
  \textbf{Pick.} For each of the \(l_k\) shrunk vectors, find the
  nearest of the \(V\) codebook entries \(E_k[1..V]\); its index is the
  code. \(c_k\) is a list of \(l_k\) integers, for example \((5, 12)\).
\item
  \textbf{Expand and refine.} Look the codes back up, stretch them to
  \(L\) positions by linear interpolation (\(U_k\)), and apply
  \(\phi_k(h) = (1-\rho)\,h + \rho\,\mathrm{conv}_k(h)\): a small
  learned convolution blended half and half (\(\rho = 0.5\)) with its
  input, which smooths the blocky stretched result.
\item
  \textbf{Subtract.} Remove what this level explained, so the next level
  only describes what is left. \(\mathrm{sg}\) (stop-gradient) means the
  subtraction is not differentiated.
\end{enumerate}

The reconstruction after level \(k\) is the sum of everything explained
so far:

\[\hat z_k = \sum_{j \le k} q_j, \qquad \hat z = \hat z_{K-1} \approx z .\]

Decoding \(\hat z_k\) for \(k < K-1\) gives the ``after level \(k\)''
texts of Figure 1.

\textbf{Codebook updates.} Codebooks are not trained by gradients. After
each step every entry moves towards the average of the vectors assigned
to it:

\[N_j \leftarrow \gamma N_j + (1-\gamma)\, n_j, \qquad m_j \leftarrow \gamma m_j + (1-\gamma) \sum_{i:\,c_i = j} x_i, \qquad E[j] = m_j / \tilde N_j .\]

\emph{Reading it:} \(n_j\) counts how many vectors picked entry \(j\) in
this batch, \(N_j\) is its running average with decay \(\gamma = 0.99\),
\(m_j\) the running sum of those vectors, and \(\tilde N_j\) is \(N_j\)
with a tiny amount added to every count (Laplace smoothing) so nothing
is divided by zero. Entries nobody uses (\(N_j\) below a threshold)
restart at a random input vector.

\textbf{Losses.} The decoder, one transformer block, reads
\(\tilde z = z + \mathrm{sg}(\hat z - z)\). Numerically this \emph{is}
\(\hat z\), but its gradient flows into \(z\) as if quantisation were
the identity (the straight-through estimator), so the encoder learns.
The training loss is

\[\mathcal{L}_{\mathrm{codec}} = \underbrace{\mathrm{CE}\big(\mathrm{Dec}(\tilde z), x\big)}_{\text{rebuild the tokens}} + \underbrace{\frac{1}{K}\sum_k \Big(\beta\,\lVert \mathrm{sg}(\hat z_k) - z\rVert^2 + \lVert \hat z_k - \mathrm{sg}(z) \rVert^2\Big)}_{\text{VQ loss}}, \quad \beta = 0.25 .\]

\emph{Reading it:} the first term asks the decoder to rebuild the exact
tokens. The \emph{commitment} term (\(\beta\)) pulls the encoder output
towards what the codes can express; the \emph{fit} term trains the
refinement convolutions \(\phi_k\) so each partial sum \(\hat z_k\) is
as close to \(z\) as possible.

\textbf{Fine-level dropout.} Half of the time the decoder receives
\(\hat z_k\) for a random \(k < K-1\) instead of the full \(\hat z\). It
must then make sense of a coarse-only description, which is why every
prefix of levels decodes to a plausible sequence: a form of nested
dropout (\citeproc{ref-nested2014}{Rippel et al. 2014}).

\textbf{Learned level 0.} Averaging an entire Sudoku grid into one
vector gives nearly the same vector for every grid, so level 0 carries
no information (Section 7.2). Upstream replaces \(D_0\) and \(U_0\) by a
strided convolution and a transposed convolution initialised to behave
like pooling and repetition. Picking a code has no gradient, so the
gradient reaches them through a straight-through estimator: \(q_0\) is
built from \(x_0 + \mathrm{sg}(E_0[c_0] - x_0)\), whose value is the
code vector and whose gradient is that of \(x_0\).

\subsection{Generator: predicting one level from the levels
above}\label{generator-predicting-one-level-from-the-levels-above}

The generator is a transformer that reads one long row of vectors:

\[[\;\text{prefix } (P)\;|\;\text{level-0 queries } (l_0)\;|\;\text{input}_1\, (l_1)\;|\;\dots\;|\;\text{input}_{K-1}\,(l_{K-1})\;],
\qquad \text{input}_k = D_k\, \hat z_{k-1} .\]

\emph{Reading it:} the \emph{prefix} holds the conditioning, for Sudoku
the 81 cells of the puzzle (blanks as 0). Level 0 has nothing above it,
so it reads learned query vectors. Every other level reads the
reconstruction from all coarser levels, shrunk to its own length. Toy
example: \(\mathrm{input}_2\) is \(\hat z_1\) (levels 0 and 1 combined)
averaged down to 4 positions. Head \(k\) turns each of the \(l_k\)
positions of its block into a probability for each of the \(V\) codes of
level \(k\).

\textbf{Training loss.} With the true codes \(c_k\) of a training
example,

\[\mathcal{L}_{\mathrm{gen}} = \frac{\sum_k l_k \cdot \mathrm{CE}_k}{\sum_k l_k}, \qquad \mathrm{CE}_k = -\frac{1}{l_k}\sum_{i=1}^{l_k} \log p_\theta\big(c_{k,i} \mid \text{input}_k, \text{prefix}\big).\]

\emph{Reading it:} \(\mathrm{CE}_k\) is the average surprise at the true
codes of level \(k\); weighting by \(l_k\) makes every code count
equally (the paper's \(\alpha = 0\) setting
(\citeproc{ref-ncm2026}{Kumbong et al. 2026}, App. F.2.3)). A uniform
guess among 16 codes scores \(\ln 16 \approx 2.77\).

\textbf{Attention pattern.} Positions of level \(k\) may read the prefix
and the other positions of level \(k\), nothing else
(\emph{block-diagonal} attention); coarser levels reach level \(k\) only
through \(\mathrm{input}_k\). Two consequences follow. Training on the
whole row at once computes exactly what level-by-level generation
computes (tested, Section 6.2). And each level's codes are chosen
without seeing each other's choices, the source of both the speed and
errors like ``red'' → ``ret''. RoPE (\citeproc{ref-rope2021}{Su et al.
2021}) positions restart at zero in every level, a learned level
embedding tells levels apart, and queries and keys are RMS-normalised.

\textbf{Conditioning corruption.} In training the generator reads true
coarser codes; in generation it reads its own, sometimes wrong, codes.
Upstream narrows this gap by replacing a fraction \(p_k\) of the
conditioning codes of level \(k\) with random codes of the same level,
keeping the true codes as targets, with \(p_k\) falling linearly across
levels (\citeproc{ref-ncm2026}{Kumbong et al. 2026}, Sec. 3.4):

\[p_k = p_{\mathrm{first}} + (p_{\mathrm{last}} - p_{\mathrm{first}})\,\frac{k}{K-2}, \qquad k = 0, \dots, K-2 .\]

This is the coarse-to-fine analogue of scheduled sampling
(\citeproc{ref-ss2015}{Bengio et al. 2015}). The paper notes that
corruption ``based on the model's actual errors may better match
inference conditions'' but does not test it; experiment E9 tests it and
finds no gain at our scale (Section 8).

\subsection{Generation}\label{generation}

For \(k = 0, \dots, K-1\): run the generator on \(\mathrm{input}_k\) and
the prefix, sample all \(l_k\) codes of level \(k\) at once, add their
contribution \(q_k\) to the reconstruction and shrink it into
\(\mathrm{input}_{k+1}\). Finally decode \(\hat z_{K-1}\) to tokens.
Cost: \(K\) generator passes and one decoder pass, independent of how
long each level is.

\begin{longtable}[]{@{}
  >{\raggedright\arraybackslash}p{(\linewidth - 4\tabcolsep) * \real{0.3333}}
  >{\raggedright\arraybackslash}p{(\linewidth - 4\tabcolsep) * \real{0.3333}}
  >{\raggedright\arraybackslash}p{(\linewidth - 4\tabcolsep) * \real{0.3333}}@{}}
\caption{Initial E6 settings versus released Original NCM Sudoku
settings (\citeproc{ref-ncm2026}{Kumbong et al. 2026}, Tables 5--6).
Later runs change exposure and corruption; see Appendix
A.}\tabularnewline
\toprule\noalign{}
\begin{minipage}[b]{\linewidth}\raggedright
Component
\end{minipage} & \begin{minipage}[b]{\linewidth}\raggedright
Original NCM (Kumbong et al.)
\end{minipage} & \begin{minipage}[b]{\linewidth}\raggedright
MLX starting recipe (E6)
\end{minipage} \\
\midrule\noalign{}
\endfirsthead
\toprule\noalign{}
\begin{minipage}[b]{\linewidth}\raggedright
Component
\end{minipage} & \begin{minipage}[b]{\linewidth}\raggedright
Original NCM (Kumbong et al.)
\end{minipage} & \begin{minipage}[b]{\linewidth}\raggedright
MLX starting recipe (E6)
\end{minipage} \\
\midrule\noalign{}
\endhead
\bottomrule\noalign{}
\endlastfoot
Tokens per board & 89 (cells and separators) & 81 (cells) \\
Levels × codes & 12 × 16, \(l_{K-1} = 89\) & 12 × 16,
\(l_{K-1} = 81\) \\
Level-0 sampler & learned (one layer, as 89 is prime) & learned from E6,
one layer \\
Codec training & 10,000 updates × batch 1,024, lr 5e-4 & 3,000 × 128, lr
2e-4 \\
Generator & 8 × 512 (25.5M), 50,000 updates × 1,024, lr 3e-4 & 8 × 256
(6.4M), 4,000 × 128, lr 6e-4 \\
Conditioning corruption & yes, linear schedule & from E8 (0.30 →
0.05) \\
Models & one per difficulty & one for all difficulties \\
\end{longtable}

\section{Implementation on Apple
Silicon}\label{implementation-on-apple-silicon}

\begin{quote}
\textbf{In short.} Pure MLX and NumPy, no code copied from upstream. Two
findings matter beyond this project: Apple GPUs compute fp32 matrix
products at reduced precision, and block-diagonal attention is faster
computed per level than through a mask.
\end{quote}

\subsection{Framework mapping}\label{framework-mapping}

Codebooks and their running averages are frozen MLX arrays updated
outside automatic differentiation; the shrink and stretch operators
\(D_k\), \(U_k\) are precomputed matrices; attention uses MLX's fused
\texttt{scaled\_dot\_product\_attention}. Weights are stored as
safetensors. The correspondence to upstream is documented component by
component in \texttt{docs/research/upstream-map.md}.

\subsection{Exactness and Metal
precision}\label{exactness-and-metal-precision}

A test asserts that level-by-level generation produces the same logits
as the training forward. On the Metal GPU the two disagreed by about
\(5 \times 10^{-4}\). Against a float64 NumPy reference, the same fp32
computation had an error of \(1.4 \times 10^{-2}\) on the GPU and
\(2 \times 10^{-6}\) on the CPU: fp32 matrix products on Metal run at
reduced internal precision, about \(10^{-3}\) relative. Exactness tests
therefore run on the CPU, and a separate test checks the GPU path within
\(5 \times 10^{-3}\). This matters for anyone writing equivalence tests
for MLX models; sampling quality is evaluated separately, and these
exactness checks do not establish the quality effect of precision.

\subsection{Segmented block-diagonal
attention}\label{segmented-block-diagonal-attention}

With a mask, attention over the full Sudoku row forms a
\(395 \times 395\) score matrix per head, most of it masked out.
Computing attention separately for the prefix and for each level (keys
and values: prefix plus that level) gives the same result with
\(P^2 + \sum_k l_k (P + l_k) \approx 48{,}000\) scores instead of
\(156{,}000\). In training (width 256, batch 128, compiled) step time
fell from 1.59 s to 1.18 s and peak memory from 18.7 GB to 16.4 GB,
measured back to back. A test checks equality with the masked
computation.

\subsection{Compilation and precision}\label{compilation-and-precision}

Compiling the training step with \texttt{mx.compile} cut step time by
about 1.3× (1.22 s → 0.92 s at width 256, batch 128). Casting the
generator to bfloat16 gave a further 1.4--1.8× in a first test but saved
no memory, because the frozen fp32 codec promotes the generator's inputs
back to fp32; mixed precision is not used in the experiments reported
here.

\subsection{Interrupted training and
reproducibility}\label{interrupted-training-and-reproducibility}

Long runs need recoverable training state as well as model weights.
Codec and generator snapshots now store optimizer moments and step,
random state, shuffled-batch position, elapsed training time and model
parameters; codec snapshots include VQ running averages and generator
snapshots retain the best-validation model. Files are replaced
atomically, with a preceding snapshot retained. A pause signal saves
after the current step; a hard crash recovers from the last completed
snapshot. Short CPU tests give identical saved state for uninterrupted
and resumed training, including the compiled generator path. This does
not establish equality across hardware or backends. Rented-pod recovery
also depends on copying snapshots outside the container before it
disappears.

E17 exposed the distinction: closing its terminal killed training after
the last saved step, 30,000 of 32,000. That legacy checkpoint contained
weights alone. Its continuation retains the learning-rate and batch
positions but resets optimizer moments and random state, so it must be
reported as a weights-only recovery. The recovered model solves 71.25 /
31.75 / 3.70\% of easy/medium/hard boards, close to E13, after 2,000
continuation steps taking 1,247.53 s. This is consistent with puzzle
construction affecting transfer, but the reset prevents a controlled
causal conclusion. The operational procedure is documented in
\texttt{docs/RECOVERY.md}.

\section{Experiments}\label{experiments}

\begin{quote}
\textbf{In short.} On a small invented language NCM is faster and
somewhat less accurate than an AR model with similar generator size. On
Sudoku our AR result is near the Original AR baseline; our NCM does not
match Original NCM's reported result.
\end{quote}

Runs use an Apple M5 Pro (64 GB unified memory, macOS 26.5), an RTX 3080
or a rented A40, as labelled for each experiment, with MLX 0.32.3. Other
GPU workloads often ran on the Mac; timings this affects are marked.

\subsection{TinyWorld: an invented
language}\label{tinyworld-an-invented-language}

TinyWorld has 15,360 sentences of the form \emph{name verb count
adjective noun}, with number agreement (``one red cat'', ``two red
cats''), at most 32 characters. 12,000 sentences train the models; the
rest are held out. A sentence is \emph{valid} if it parses and agrees.
The codec (6 levels, \(l = 1, 2, 4, 8, 16, 32\), 128 codes each, 1.3M
parameters) rebuilds 100\% of held-out sentences exactly. The NCM
generator (5.0M parameters) is compared with a character-level AR
transformer of the same width and depth (4.8M parameters) with a KV
cache.

\begin{longtable}[]{@{}
  >{\raggedright\arraybackslash}p{(\linewidth - 4\tabcolsep) * \real{0.3333}}
  >{\raggedright\arraybackslash}p{(\linewidth - 4\tabcolsep) * \real{0.3333}}
  >{\raggedright\arraybackslash}p{(\linewidth - 4\tabcolsep) * \real{0.3333}}@{}}
\caption{TinyWorld, NCM and AR with similar generator size (E1); NCM
also retains its codec.}\tabularnewline
\toprule\noalign{}
\begin{minipage}[b]{\linewidth}\raggedright
\end{minipage} & \begin{minipage}[b]{\linewidth}\raggedright
MLX implementation
\end{minipage} & \begin{minipage}[b]{\linewidth}\raggedright
MLX AR baseline
\end{minipage} \\
\midrule\noalign{}
\endfirsthead
\toprule\noalign{}
\begin{minipage}[b]{\linewidth}\raggedright
\end{minipage} & \begin{minipage}[b]{\linewidth}\raggedright
MLX implementation
\end{minipage} & \begin{minipage}[b]{\linewidth}\raggedright
MLX AR baseline
\end{minipage} \\
\midrule\noalign{}
\endhead
\bottomrule\noalign{}
\endlastfoot
Valid, temperature 1.0 & 89.8\% & 98.6\% \\
Valid, temperature 0.7 & 95.8\% & 99.8\% \\
Valid and not in the training set, temperature 1.0 & 16.3\% & 23.5\% \\
Forward passes per sequence & 6 + 1 decode & 32 \\
Median latency, batch 1 / 16 / 256 & 4.5 / 5.8 / 37.8 ms & 30.5 / 22.9 /
49.8 ms \\
\end{longtable}

\begin{figure}
\centering
\includegraphics[width=1\linewidth,height=\textheight,keepaspectratio,alt={Median
latency per batch for NCM and AR models. Left: TinyWorld. Right: Sudoku.
Both axes logarithmic.}]{figures/latency.pdf}
\caption{Median latency per batch for NCM and AR models. Left:
TinyWorld. Right: Sudoku. Both axes logarithmic.}
\end{figure}

NCM is less accurate and faster. Its errors are local spelling slips;
agreement, a sentence-wide rule, holds in 99.8\% of well-formed samples.
Forcing the true codes for the first \(k\) levels and sampling the rest
(E2) shows that errors arise at every level, not only the finest.
Letting level \(k\) also read coarser levels (block-causal attention,
E3) raised validity at temperature 1 from 89.8\% to 92.2\% (n = 1,024,
standard error about 0.9 points) at 1.5--1.9× the latency.

\subsection{Sudoku}\label{sudoku}

\textbf{Comparison labels.} ``MLX implementation'' means the
independently trained NCM model in this repository. ``Original NCM''
means Kumbong et al.'s published model and checkpoint results
(\citeproc{ref-ncm2026}{Kumbong et al. 2026}). The ``MLX AR baseline''
is trained here; the ``Original AR baseline'' is quoted from their
study. The underlying method is NCM in both implementations.

Each board is 81 digits; the puzzle (blanks as 0) is the prefix and the
solution is the target. Training uses 200,000 synthetic puzzles from
randomised backtracking with 28--45 givens. Test sets are the 2,000
easy, medium and hard boards (40, 35 and 30 givens) published with the
NCM paper. A board counts as \emph{solved} if the output is a valid grid
that keeps every given. Decoding is greedy.

\begin{longtable}[]{@{}
  >{\raggedright\arraybackslash}p{(\linewidth - 8\tabcolsep) * \real{0.2000}}
  >{\raggedright\arraybackslash}p{(\linewidth - 8\tabcolsep) * \real{0.2000}}
  >{\raggedright\arraybackslash}p{(\linewidth - 8\tabcolsep) * \real{0.2000}}
  >{\raggedright\arraybackslash}p{(\linewidth - 8\tabcolsep) * \real{0.2000}}
  >{\raggedright\arraybackslash}p{(\linewidth - 8\tabcolsep) * \real{0.2000}}@{}}
\caption{Sudoku, share of boards solved. Params counts generators only
for NCM; the codec is additional. Paper rows are reference results from
(\citeproc{ref-ncm2026}{Kumbong et al. 2026}, Table 3). Budgets differ
except for the explicitly matched E11 comparison.}\tabularnewline
\toprule\noalign{}
\begin{minipage}[b]{\linewidth}\raggedright
Model
\end{minipage} & \begin{minipage}[b]{\linewidth}\raggedright
Params
\end{minipage} & \begin{minipage}[b]{\linewidth}\raggedright
Easy
\end{minipage} & \begin{minipage}[b]{\linewidth}\raggedright
Medium
\end{minipage} & \begin{minipage}[b]{\linewidth}\raggedright
Hard
\end{minipage} \\
\midrule\noalign{}
\endfirsthead
\toprule\noalign{}
\begin{minipage}[b]{\linewidth}\raggedright
Model
\end{minipage} & \begin{minipage}[b]{\linewidth}\raggedright
Params
\end{minipage} & \begin{minipage}[b]{\linewidth}\raggedright
Easy
\end{minipage} & \begin{minipage}[b]{\linewidth}\raggedright
Medium
\end{minipage} & \begin{minipage}[b]{\linewidth}\raggedright
Hard
\end{minipage} \\
\midrule\noalign{}
\endhead
\bottomrule\noalign{}
\endlastfoot
MLX implementation, E4--E6 & 6.4M & 0.0\% & 0.0\% & 0.0\% \\
MLX implementation, E8 (+ corruption) & 6.4M & 6.2\% & 0.35\% & 0.0\% \\
MLX implementation, E11 (+ 4× data) & 6.4M & 55.2\% & 18.1\% & 1.6\% \\
MLX implementation, E13 (+ stronger corruption) & 6.4M & 70.9\% & 31.2\%
& 4.0\% \\
MLX implementation, E14 (+ 16× data) & 6.4M & 87.4\% & 53.4\% &
11.2\% \\
MLX implementation, E15 & 6.4M & 96.25\% & 76.65\% & 31.90\% \\
MLX implementation, E18 (base data) & 6.4M & 84.25\% & 49.10\% &
10.55\% \\
MLX implementation, E19 pilot & 6.4M & 90.20\% & 59.75\% & 16.55\% \\
MLX implementation, E19 funded (seed 0) & 6.4M & 98.10\% & 83.40\% &
40.30\% \\
MLX implementation, E20 (seed 1) & 6.4M & 97.25\% & 79.60\% & 33.25\% \\
MLX AR baseline & 6.3M & 13.8\% & 3.0\% & 0.4\% \\
MLX AR baseline, E11's budget (2M sequences) & 6.3M & 20.5\% & 4.9\% &
0.5\% \\
Qwen3.8 27B, 4-bit, no reasoning (50 boards each) & 27B & 0\% & 0\% &
0\% \\
Original AR baseline (Kumbong et al.) & 25M & 14.6\% & 5.1\% & 1.0\% \\
MDLM (reported by Kumbong et al.) & 29M & 92.0\% & 77.1\% & 30.2\% \\
Original NCM (Kumbong et al.) & 25M & 99.5\% & 96.5\% & 86.1\% \\
\end{longtable}

The first AR row is within a few percentage points of the Original AR
baseline, but this numerical proximity does not establish a reproduction
of its model or training recipe. Without conditioning corruption
(E4--E6) our NCM solves nothing; with it (E8) it solves 6.2\% of easy
boards and fills 74.7\% of blank cells correctly, more than the AR
model's 71.0\%, while still solving fewer boards. E11 solves 55.2\% of
easy boards; at the same 2.048M generator presentations, the AR model
solves 20.5\%, 4.9\% and 0.5\%. This matches generator data exposure,
not total training cost, FLOPs or system size. Section 8 records the
changes.

The shared E6 codec contains 13,816,704 gradient-trained values, 49,152
EMA-trained codebook values and 49,344 EMA accumulators: 13,915,200
saved scalar values in total. Its learned first-level down/up
convolutions account for 10,617,344 trainable values. The 6,440,128
generator values are additional; generator-size comparisons exclude the
codec. Older rows preserve historical rounding; E15 and the newly
completed runs retain their two-decimal solved rates. Each test set has
2,000 boards. Most recipes use one training seed; E19-funded and E20
provide two seeds of the same scientific recipe. Repeated sampling seeds
do not quantify training variability.

Historical synchronized timings use temperature 1.0 while Sudoku quality
uses greedy temperature 0.0. Both samplers remain uncompiled;
\texttt{mx.compile} is used only for generator training when configured.
Latency therefore does not directly time the solved-board decoding
setting. The benchmark includes the output host copy, prefix work and
final decoder, and records shared-device contention; cross-hardware
timing ratios are not controlled comparisons.

\textbf{Smaller-generator result.} Easy-board accuracy is within 1.40
percentage points of Original NCM (98.10\% versus 99.50\%) with a 6.44M
versus 25.46M generator: 3.95× smaller and 3.125× fewer generator
presentations. This is a useful small-model result, not overall task
parity: medium and hard remain 13.10 and 45.80 points behind. The
separate codec is excluded from this size ratio, and the recipes and
representation differ. E20 retains the same generator size and exposure
but its easy score is 2.25 points below Original NCM; its medium/hard
gaps are 16.90/52.85 points. The best-seed proximity on easy boards is
not a reproducibility guarantee.

\subsubsection{Completed E18 and E19
experiments}\label{completed-e18-and-e19-experiments}

E18 tests stronger corruption on the base 200,000-puzzle corpus at E15's
8.192M-presentation budget. Its 84.25 / 49.10 / 10.55\% solve rates
remain below both E14 and E15. Stronger corruption is therefore not a
universal improvement, and this negative result remains part of the
evidence.

The short E19 pilot uses the frozen E6 codec and 8M cached unique
puzzles in addition to the base corpus. It completes 23,895 updates at
batch 128 (3,058,560 presentations), solving 90.20 / 59.75 / 16.55\%.
E19-funded starts a fresh generator with the same corpus and a separate
128,000-update cosine schedule. At 16,384,000 presentations it solves
98.10 / 83.40 / 40.30\%, improving on E15 by 1.85 / 6.75 / 8.40
percentage points. The exposure doubles while batch size, corpus and
schedule differ. These comparisons establish outcomes of whole recipes,
not a controlled effect of any one change.

\begin{longtable}[]{@{}lrrrr@{}}
\caption{Completed training and checkpoint selection (training-loop
seconds).}\tabularnewline
\toprule\noalign{}
Run & Updates × batch & Presentations & Best step & Train seconds \\
\midrule\noalign{}
\endfirsthead
\toprule\noalign{}
Run & Updates × batch & Presentations & Best step & Train seconds \\
\midrule\noalign{}
\endhead
\bottomrule\noalign{}
\endlastfoot
E18 & 128,000 × 64 & 8,192,000 & 128,000 & 23,366.69 \\
E19 pilot & 23,895 × 128 & 3,058,560 & 23,895 & 10,225.84 \\
E19 funded & 128,000 × 128 & 16,384,000 & 127,000 & 54,419.81 \\
\end{longtable}

The funded run completes its full schedule without a time-limit stop or
optimizer reset; its final resumable snapshot is step 128,000, while
held-out evaluation uses the best-validation step 127,000. Best weighted
validation CE is 0.133244. The complete remote pipeline takes 55,700.91
s, including preprocessing, evaluation, benchmarking and transfer but
excluding rental startup and the final cleanup wait. All 17 downloaded
report/artifact files match their completion hashes. Source revision is
\texttt{80dd445e5d17f54f53687a88459e011de4f9006f}; the generator SHA-256
and full result extract are retained in the repository's E19 report.
Full-state snapshots preserve optimizer and random state; this
uninterrupted run is not a test of cross-backend exact recovery.

Funded validation solves 70.90\% of 1,000 boards. On easy/medium/hard
tests, blank-cell accuracy is 99.69 / 97.12 / 86.99\% and givens
preservation is 100 / 100 / 99.60\%; full-depth codec validation
reconstruction is separately 100\%. Wilson 95\% intervals for the solved
rates are 97.40--98.61 / 81.71--84.97 / 38.17--42.47\%. These assume
independent test boards, exclude training-seed variability and do not
establish paired significance against prior runs. Per-board prediction
archives were not retained. E20 supplies a second training seed below.

\subsubsection{E20: a second training
seed}\label{e20-a-second-training-seed}

E20 repeats the E19-funded scientific recipe with configured training
seed 1 instead of 0. Both complete 128,000 updates at batch 128, using
the frozen E6 codec, 200,000 base puzzles plus 8M cached unique puzzles,
corruption 0.50→0.10, learning rate 0.0006 and 1,000-step warmup. Shared
codec and corpus input-manifest hashes match. The configured seed
affects initialization, minibatch order, dropout and corruption;
internal generator streams use that seed plus one. These influences are
not separately controlled.

The run name, output path and operational cap also change. E20's cap is
960 rather than 1,440 minutes; neither run reaches its cap. E20 records
no optimizer reset or truncation, and its final full-state snapshot and
selected best checkpoint are both step 128,000. Training-loop time is
54,122.43 s; pipeline time is 55,543.62 s under the E19 timing boundary.
Best weighted code CE is 0.142558, compared with 0.133244 for seed 0.
Separate physical A40 allocations and missing remote source inventories
prevent a strict same-host causal attribution.

\begin{longtable}[]{@{}
  >{\raggedright\arraybackslash}p{(\linewidth - 6\tabcolsep) * \real{0.2000}}
  >{\raggedleft\arraybackslash}p{(\linewidth - 6\tabcolsep) * \real{0.2667}}
  >{\raggedleft\arraybackslash}p{(\linewidth - 6\tabcolsep) * \real{0.2667}}
  >{\raggedleft\arraybackslash}p{(\linewidth - 6\tabcolsep) * \real{0.2667}}@{}}
\caption{Two-seed held-out aggregates. Each seed reports the same
2,000-board benchmark per difficulty; the mean is descriptive, not a
pooled independent sample.}\tabularnewline
\toprule\noalign{}
\begin{minipage}[b]{\linewidth}\raggedright
Quantity
\end{minipage} & \begin{minipage}[b]{\linewidth}\raggedleft
Easy
\end{minipage} & \begin{minipage}[b]{\linewidth}\raggedleft
Medium
\end{minipage} & \begin{minipage}[b]{\linewidth}\raggedleft
Hard
\end{minipage} \\
\midrule\noalign{}
\endfirsthead
\toprule\noalign{}
\begin{minipage}[b]{\linewidth}\raggedright
Quantity
\end{minipage} & \begin{minipage}[b]{\linewidth}\raggedleft
Easy
\end{minipage} & \begin{minipage}[b]{\linewidth}\raggedleft
Medium
\end{minipage} & \begin{minipage}[b]{\linewidth}\raggedleft
Hard
\end{minipage} \\
\midrule\noalign{}
\endhead
\bottomrule\noalign{}
\endlastfoot
Seed 0 solved boards & 1,962 & 1,668 & 806 \\
Seed 1 solved boards & 1,945 & 1,592 & 665 \\
Seed 1 solved rate & 97.25\% & 79.60\% & 33.25\% \\
Seed 1 minus seed 0 & −0.85 pp & −3.80 pp & −7.05 pp \\
Two-seed mean & 97.675\% & 81.500\% & 36.775\% \\
Seed 1 Wilson 95\% interval & 96.44--97.88\% & 77.78--81.31\% &
31.22--35.35\% \\
\end{longtable}

Both seeds retain strong easy-board performance and substantially weaker
hard performance. E19-funded remains the best observed seed; neither
two-seed means nor board-level Wilson intervals reliably estimate the
population distribution of training outcomes. The intervals also exclude
related-board dependence and repeated benchmark-guided recipe selection.
Without per-board predictions, paired significance and agreement between
the solved-board sets cannot be computed. No additional model evaluation
is included in this version.

E20 validation solves 71.00\%, compared with 70.90\% for E19-funded,
while hard blank-cell accuracy falls from 86.99\% to 84.92\%. Similar
validation solving therefore does not imply similar external hard-set
performance. The saved hard aggregates identify 670 valid grids, of
which 665 preserve the givens: 1,330 failures are invalid grids and five
additional failures are valid grids contradicting clues. Fourteen boards
violate givens overall. Clue preservation alone cannot repair the
dominant invalid-grid failure; the aggregates do not identify its
internal cause.

\textbf{Audit boundary.} All 15 retained E20 completion files match
recorded hashes; inference tensors equal the saved best-model tensors,
and optimizer/training steps agree at 128,000. Generator SHA-256 is
\texttt{b6ddadb7278e7f1c5542eecef676840f39cb829e802d189b65f2dadd645c2b3d};
the declared source is
\texttt{4ad3646189197fe1ca76e45308848793fefe9a3f}. These checks
establish retained-byte integrity and internal consistency, not
independent reproduction. No checkpoint-bound ordered scoring manifest,
per-board predictions or remote source digest inventory was retained.
Current local board files match the published hashes, but they were
absent from the remote transfer-verification manifest. Actual remote
evaluation identity, complete duplicate-free coverage and exact source
bytes remain independently unverified. Reported outcomes are historical
aggregates with these limits.

E20's recorded synchronized benchmark uses temperature 1, three warmups
and 20 iterations. Batch 1/16/64/256 p50 is
29.453/48.506/174.326/581.791 ms; p95 is 30.323/49.274/175.352/583.397
ms. The CUDA GPU-load log is empty because the wrapper used a Metal-only
probe. Contention at benchmark time is unknown, and these timings
neither measure greedy quality generation nor establish a seed-dependent
speed effect.

\begin{figure}
\centering
\includegraphics[width=1\linewidth,height=\textheight,keepaspectratio,alt={Completed
MLX implementation recipes, with rounded generator presentations beneath
each run; exact exposures are in the tables. Each bar reports 2,000
held-out boards; retained aggregate and provenance limitations apply.
Dashed lines show Original NCM (Kumbong et al.), with its larger
generator and 51.2M presentations; they are not matched-budget
targets.}]{figures/sudoku-completed-runs.pdf}
\caption{Completed MLX implementation recipes, with rounded generator
presentations beneath each run; exact exposures are in the tables. Each
bar reports 2,000 held-out boards; retained aggregate and provenance
limitations apply. Dashed lines show Original NCM (Kumbong et al.), with
its larger generator and 51.2M presentations; they are not
matched-budget targets.}
\end{figure}

Funded synchronized generation at temperature 1.0 uses three warmups and
20 timed iterations, 12 generator forwards and completion through a host
copy. Batch 1/16/64/256 p50 latency is 26.477 / 47.905 / 174.279 /
583.670 ms; p95 is 26.690 / 48.308 / 174.552 / 584.715 ms. This timing
does not measure greedy quality decoding. Preflight records zero GPU
utilization, zero allocated memory and no compute applications; an
immediate pre-benchmark contention sample is not retained, so these are
not certified idle-device limits. Neither cross-host A40 timings nor RTX
3080 comparisons isolate a hardware speed ratio.

The first funded allocation fails during shell preflight, before any
training step, and is deleted. The failure is SIGPIPE from an
early-closing version-check pipeline under \texttt{pipefail}; three
regression cases reproduce it, and the corrected fresh attempt passes
all 78 tests. Image startup timeouts recover before readiness. These are
setup incidents, not extra training runs or reset-state continuations.

\begin{figure}
\centering
\includegraphics[width=0.62\linewidth,height=\textheight,keepaspectratio,alt={Cell
accuracy when decoding from the first k levels of the codec, against the
share of codes used. E4 stores almost all information in one level; E5
and E6 refine gradually.}]{figures/codec-refinement.pdf}
\caption{Cell accuracy when decoding from the first \(k\) levels of the
codec, against the share of codes used. E4 stores almost all information
in one level; E5 and E6 refine gradually.}
\end{figure}

\begin{figure}
\centering
\includegraphics[width=0.62\linewidth,height=\textheight,keepaspectratio,alt={Validation
cross-entropy per code during generator
training.}]{figures/sudoku-val-curves.pdf}
\caption{Validation cross-entropy per code during generator training.}
\end{figure}

\begin{figure}
\centering
\includegraphics[width=1\linewidth,height=\textheight,keepaspectratio,alt={Blank
cells filled with the reference digit, per test set. Chance is about
11\%.}]{figures/sudoku-blank-cells.pdf}
\caption{Blank cells filled with the reference digit, per test set.
Chance is about 11\%.}
\end{figure}

\textbf{Why per-cell gains do not become solved boards.} A board is
about 300 codes. Teacher-forced accuracy per level in E6 ranges from
25\% (level 0, which the puzzle barely determines) to 99.6\% (the last
level); generating greedily, an error at an early level changes the
inputs of every later level. E8's corruption training attacks exactly
this compounding. Even 98\% accuracy per cell, with independent errors,
would solve only \(0.98^{41} \approx 44\%\) of easy boards, which have
41 blanks. The paper's generator was trained on about 51 million
sequences; E6 saw 0.5 million.

\subsection{A large chat model as
reference}\label{a-large-chat-model-as-reference}

Qwen3.8 27B (4-bit, served locally) answered each board in one request
at temperature 0 with reasoning disabled. It solved none of 150 boards;
its blank-cell accuracy was 11--15\%, near chance, and about half its
replies were not a 9 × 9 grid. With reasoning enabled and up to 6,000
tokens per board it did not finish ten boards in 95 minutes on the
shared GPU. This is a reference point, not a matched comparison: the
model is about 4,000 times larger and was not trained on the task.

\section{The journey: experiment by
experiment}\label{the-journey-experiment-by-experiment}

\begin{quote}
\textbf{In short.} The experiments examine architecture, training and
data choices; some combine multiple changes. Experiment numbers (E0, E1,
\ldots) identify runs; the full record with every number is in
\texttt{docs/EXPERIMENTS.md}.
\end{quote}

\begin{longtable}[]{@{}
  >{\raggedright\arraybackslash}p{(\linewidth - 8\tabcolsep) * \real{0.2000}}
  >{\raggedright\arraybackslash}p{(\linewidth - 8\tabcolsep) * \real{0.2000}}
  >{\raggedright\arraybackslash}p{(\linewidth - 8\tabcolsep) * \real{0.2000}}
  >{\raggedright\arraybackslash}p{(\linewidth - 8\tabcolsep) * \real{0.2000}}
  >{\raggedright\arraybackslash}p{(\linewidth - 8\tabcolsep) * \real{0.2000}}@{}}
\caption{Experiments in order. Sudoku results are for the easy test
boards.}\tabularnewline
\toprule\noalign{}
\begin{minipage}[b]{\linewidth}\raggedright
Run
\end{minipage} & \begin{minipage}[b]{\linewidth}\raggedright
Question
\end{minipage} & \begin{minipage}[b]{\linewidth}\raggedright
Change from the previous run
\end{minipage} & \begin{minipage}[b]{\linewidth}\raggedright
Result
\end{minipage} & \begin{minipage}[b]{\linewidth}\raggedright
What we learned
\end{minipage} \\
\midrule\noalign{}
\endfirsthead
\toprule\noalign{}
\begin{minipage}[b]{\linewidth}\raggedright
Run
\end{minipage} & \begin{minipage}[b]{\linewidth}\raggedright
Question
\end{minipage} & \begin{minipage}[b]{\linewidth}\raggedright
Change from the previous run
\end{minipage} & \begin{minipage}[b]{\linewidth}\raggedright
Result
\end{minipage} & \begin{minipage}[b]{\linewidth}\raggedright
What we learned
\end{minipage} \\
\midrule\noalign{}
\endhead
\bottomrule\noalign{}
\endlastfoot
E0 & Does the pipeline work end to end? & First version, 4,000 TinyWorld
sentences & 94.6\% valid; overfits & The codec rebuilds 99.4\% of
sentences; models need more data \\
E1 & NCM vs AR, fairly & 12,000 sentences, best-validation checkpoints &
NCM 89.8\% vs AR 98.6\% valid; 4--5× faster & Speed holds on a laptop;
the accuracy gap is real \\
E2 & Which level causes errors? & Force true codes for levels \(0..k\) &
Errors at every level & Not only the finest level is at fault \\
E3 & Does reading coarser levels help? & Block-causal attention & 92.2\%
valid, 1.5--1.9× slower & Small gain; kept as an option \\
E4 & Does NCM solve Sudoku? & Puzzle as prefix, 6 levels & 0\% solved;
AR 13.8\% & One level holds all information: no real coarse-to-fine \\
E5 & Does upstream's level schedule fix it? & 12 levels × 16 codes,
lookup encoder & 0\%; blank cells 52\% (was 29\%) & Gradual refinement,
but level 0 is dead \\
E6 & Does a learned level 0 fix it? & Learned level-0 sampler,
GroupNorm, lower codec lr & 0\%; blank cells 54\% & Level 0 now used;
small gains; diverges at lr 1e-3 \\
E7 & Upstream's generator width & 8 × 512 & stopped early & Recipe
search showed corruption was missing; GPU time moved to E8 \\
E8 & Does conditioning corruption help? & Uniform corruption 0.30 → 0.05
& 6.2\% solved; blank cells 74.7\% & The largest single gain: first
solved boards; beats AR on blank cells; validation loss gets worse while
generation gets better \\
E9 & Does corruption from the model's own samples help more? &
Replacements drawn from the generator & 3.0\% solved; blank cells 74.7\%
& No: same blank cells, fewer solved boards, 1.56× training time \\
E10 & Upstream's width & 8 × 512, lr 3e-4, 1.5M sequences, on the Mac &
55.9\% / 18.4\% / 2.0\% & Width alone does not help at this budget:
matches the 256-wide E11 \\
E11 & More data & E8 on 2M sequences (4×), RTX 3080 & 55.2\% easy,
18.1\% medium, 1.6\% hard solved & Data was the largest remaining lever;
NCM now beats the same-size AR model \\
E12 & Upstream's width on a rented GPU & 8 × 512, stronger corruption,
3.25M sequences, NVIDIA A40 & 74.6\% / 35.2\% / 6.0\% & Width adds
little; data still dominates \\
E13 & Stronger corruption & 0.50 → 0.10 at E11's budget & 70.9\% /
31.2\% / 4.0\% & Corruption strength matters: +16 points on easy \\
E14 & Even more data & E11 on 8.2M sequences (16× E8) & 87.4\% / 53.4\%
/ 11.2\% & Data keeps paying; still not saturated \\
E16 & Fresh unique puzzles vs repetition & E13 on 2M greedy-unique
puzzles seen once & 57.6\% / 21.8\% / 2.2\% & Worse than E13: how
puzzles are built matters more than uniqueness \\
E17 & Unique puzzles built like the test boards & E13 on random-removal
puzzles that happen to be unique; weights-only recovery & 71.25\% /
31.75\% / 3.70\% & Close to E13; optimizer reset limits
interpretation \\
E15 & Both levers + unique-solution puzzles & E14 budget, corruption
0.50 → 0.10, +4M unique puzzles & 96.2\% / 76.6\% / 31.9\% & Within 3.3
points on easy; roughly one quarter of upstream's generator size and one
sixth of generator exposure \\
E18 & Both levers, base data only & E14 budget, corruption 0.50 → 0.10 &
84.25\% / 49.10\% / 10.55\% & Below E14 and E15; stronger corruption
alone is not sufficient \\
E19 pilot & Larger-corpus short schedule & +8M unique puzzles; 3.059M
presentations, batch 128, A40 & 90.20\% / 59.75\% / 16.55\% & Below E15
at less exposure; not a matched comparison \\
E19 funded & Larger corpus and longer schedule & Fresh initialization;
16.384M presentations, batch 128, A40 & 98.10\% / 83.40\% / 40.30\% &
Best observed seed at twice E15 exposure; still below Original NCM \\
E20 & Does the funded recipe repeat? & Fresh seed 1, same scientific
settings and exposure & 97.25\% / 79.60\% / 33.25\% & Easy/medium gains
persist; hard result varies by 7.05 points; audit limits remain \\
\end{longtable}

\textbf{What the journey shows.} The first Sudoku configuration failed
for a reason no reconstruction test would reveal: its codec rebuilds
boards perfectly yet places nearly all information in one level (the
codec-refinement figure, E4), so the generator faces one large guess
instead of a series of small ones. The fixes taken from upstream each
helped a little and none was enough on its own. Reading the paper's
appendix in full showed the two largest gaps: conditioning corruption
and a training budget about a hundred times larger. Corruption alone
(E8) produced the first solved boards and lifted blank-cell accuracy by
20 points on every test set. Its validation loss was \emph{worse} than
without corruption, a reminder that teacher-forced loss does not measure
what generation needs.

\section{Discussion}\label{discussion}

\textbf{What carries over to a laptop.} Measured Apple Silicon runs show
lower small-batch NCM latency than an AR model of similar generator size
on the tested synthetic task. The advantage shrinks with batch size and
code count; historical 12-level Sudoku generation is slower than AR from
batch 64 on. These observations do not establish natural-text speed or
isolate dispatch, memory traffic and kernel costs. The completed E19/E20
training runs use CUDA, not the Mac.

\textbf{Partial success at the tested scale.} A 6.44M generator learns
substantial puzzle-conditioned structure: both funded-recipe seeds solve
more than 97\% of easy boards and 79\% of medium boards. Both exceed
E15's reported rates, with twice its generator exposure and a changed
overall recipe. This is sufficient evidence to continue studying the
implementation. It is not an estimate of minimum required capacity, a
scaling law or a demonstration of human-style logical reasoning. The
additional codec prevents treating generator size as whole-system size.

\textbf{Why hard Sudoku remains unresolved.} The hard benchmark provides
30 clues, leaving 51 blanks, versus 40 clues and 41 blanks on easy
boards. Every constraint and input clue must be satisfied
simultaneously. Errors can be correlated, and a small change in average
cell accuracy can accompany a larger change in exact-board accuracy.
Neither clue count nor these aggregates reveals which logical
dependencies the model fails to represent. The reference comparison also
changes capacity, exposure, representation, corpus and
difficulty-specific training. Those are candidate explanations, not
experimentally established causes. Earlier width-512 trials do not
justify promising that widening alone closes the gap.

\textbf{How to challenge the interpretation.} Before further long runs,
retain checkpoint-bound input identities, per-board outputs and
per-level diagnostics. Use a fresh untouched test set after fixing the
selection protocol, since repeated examination of these benchmarks can
guide recipe choices. Additional seeds characterize variation; separate
random streams and controlled interventions are needed to identify
whether initialization, order, dropout or corruption matters. True-code
interventions on a diagnostic subset can test error propagation.
Matched-data capacity and exposure ablations can test scale without
conflating both. Checkpoint selection must use validation data rather
than the external benchmark. These are proposed tests, not completed
findings.

\textbf{Decision for the next phase.} The project can proceed to
controlled text research without declaring the Sudoku reproduction
complete. Keep the small generator as a reference while investigating
independently converted Original NCM weights. Conversion parity would
validate inference compatibility; it would not reproduce training or
validate our smaller model. Begin text with a pinned permitted corpus,
explicit tokenizer and padding boundaries, and document-separated
splits. Report reconstruction and free-running generation separately,
include readable failures and an autoregressive comparison with
total-cost qualifications. Synthetic TinyStories would be an initial
engineering benchmark; naturally occurring text and conditional tasks
require further evidence. No instruction-following, summarization,
translation or useful application integration follows from Sudoku
scores.

\section{Limitations}\label{limitations}

Most generator recipes use at most 6.4M parameters; width-512 ablations
use about 25M. E19-funded and E20 provide only two training seeds of one
recipe, not a reliable population variance estimate. The best hard score
remains 45.80 points below Original NCM; the second seed remains 52.85
points below. No current-budget AR comparison, matched total-cost study
or natural-text validation is available. Full-depth codec validation
reconstruction does not prove useful coarse representations under
generated errors. Historical aggregate-only scoring, missing remote
manifests, repeated benchmark-guided development and unknown benchmark
contention limit independent audit and generalization. The original
training implementation is not released; unstated recipe details remain
uncertain. Existing measurements establish neither production
reliability nor a benefit in a user application.

\section{Related work}\label{related-work}

Parallel and coarse-to-fine decoding for text includes discrete
diffusion (\citeproc{ref-mdlm2024}{Sahoo et al. 2024}) and the methods
compared in (\citeproc{ref-ncm2026}{Kumbong et al. 2026}). Residual
quantisation is standard in audio codecs
(\citeproc{ref-soundstream2021}{Zeghidour et al. 2021}) and image
generation (\citeproc{ref-rqvae2022}{Lee et al. 2022};
\citeproc{ref-var2024}{Tian et al. 2024}). Scheduled sampling
(\citeproc{ref-ss2015}{Bengio et al. 2015}) addresses the train--test
mismatch that conditioning corruption targets.
\href{https://github.com/Bakhaow/ncm-forge}{Bakhaow/ncm-forge} adds
PyTorch training loops around upstream modules. The repository's dated
related-work comparison records other implementations; we make no
priority claim for this MLX implementation.

\section{Conclusion}\label{conclusion}

This independent MLX implementation preserves exact CPU
training/per-level inference equivalence and demonstrates coarse-to-fine
generation on a finite synthetic language and Sudoku. The funded
small-generator recipe reports 97.25--98.10\% easy, 79.60--83.40\%
medium and 33.25--40.30\% hard solving across two seeds at equal
exposure. We regard this as partial implementation and small-model
success, with substantial failures on hard boards and unresolved
provenance and reproducibility limits. It does not reproduce Original
NCM's scores or establish matched-cost superiority. The observed seed
difference warrants better diagnostics rather than attributing the
deficit to a single component. The next research phase will test text
generation and its quality/cost tradeoff; those capabilities remain
unvalidated.

\section*{Code and reproducibility}\label{code-and-reproducibility}
\addcontentsline{toc}{section}{Code and reproducibility}

Code, configurations and this paper:
\href{https://github.com/janmus1c/ncm-mlx}{github.com/janmus1c/ncm-mlx}.
Every configuration is in \texttt{configs/}; measured runs are logged in
\texttt{docs/EXPERIMENTS.md}. The benchmark guide explains budget
matching, parameter accounting, synchronization and uncertainty.

\begin{Shaded}
\begin{Highlighting}[]
\FunctionTok{git}\NormalTok{ clone https://github.com/janmus1c/ncm{-}mlx }\KeywordTok{\&\&} \BuiltInTok{cd}\NormalTok{ ncm{-}mlx}
\ExtensionTok{mise}\NormalTok{ run setup                                     }\CommentTok{\# Python 3.12, MLX, tests}
\ExtensionTok{scripts/run\_experiment.sh}\NormalTok{ configs/sudoku{-}e6.json   }\CommentTok{\# train, evaluate, benchmark}
\ExtensionTok{.venv/bin/python}\NormalTok{ paper/make\_figures.py }\AttributeTok{{-}{-}plot}      \CommentTok{\# rebuild figures from committed data}
\ExtensionTok{paper/build.sh}                                     \CommentTok{\# rebuild this paper}
\end{Highlighting}
\end{Shaded}

\section*{AI assistance}\label{ai-assistance}
\addcontentsline{toc}{section}{AI assistance}

OpenAI Codex assisted with implementation, experiments, analysis,
figures, documentation and manuscript drafting and revision. It is not
an author. This assistance does not constitute independent verification
of the results or human review of the manuscript.

\section*{References}\label{references}
\addcontentsline{toc}{section}{References}

\protect\phantomsection\label{refs}
\begin{CSLReferences}{1}{1}
\bibitem[\citeproctext]{ref-ss2015}
Bengio, Samy, Oriol Vinyals, Navdeep Jaitly, and Noam Shazeer. 2015.
{`Scheduled Sampling for Sequence Prediction with Recurrent Neural
Networks'}. \emph{Advances in Neural Information Processing Systems}.

\bibitem[\citeproctext]{ref-mlx2023}
Hannun, Awni, Jagrit Digani, Angelos Katharopoulos, and Ronan Collobert.
2023. \emph{{MLX}: Efficient and Flexible Machine Learning on Apple
Silicon}.
\href{https://github.com/ml-explore/mlx}{Https://github.com/ml-explore/mlx}.

\bibitem[\citeproctext]{ref-ncm2026}
Kumbong, Hermann, Ishan S. Khare, Iddah Mlauzi, et al. 2026.
{`Next-Concept Modeling: Fast, Parallel Coarse-to-Fine Text
Generation'}. \emph{arXiv Preprint arXiv:2610.14162}.

\bibitem[\citeproctext]{ref-rqvae2022}
Lee, Doyup, Chiheon Kim, Saehoon Kim, Minsu Cho, and Wook-Shin Han.
2022. {`Autoregressive Image Generation Using Residual Quantization'}.
\emph{Proceedings of the IEEE/CVF Conference on Computer Vision and
Pattern Recognition}.

\bibitem[\citeproctext]{ref-vqvae2017}
Oord, Aaron van den, Oriol Vinyals, and Koray Kavukcuoglu. 2017.
{`Neural Discrete Representation Learning'}. \emph{Advances in Neural
Information Processing Systems}.

\bibitem[\citeproctext]{ref-nested2014}
Rippel, Oren, Michael Gelbart, and Ryan Adams. 2014. {`Learning Ordered
Representations with Nested Dropout'}. \emph{International Conference on
Machine Learning}.

\bibitem[\citeproctext]{ref-mdlm2024}
Sahoo, Subham Sekhar, Marianne Arriola, Yair Schiff, et al. 2024.
{`Simple and Effective Masked Diffusion Language Models'}.
\emph{Advances in Neural Information Processing Systems}.

\bibitem[\citeproctext]{ref-rope2021}
Su, Jianlin, Yu Lu, Shengfeng Pan, Ahmed Murtadha, Bo Wen, and Yunfeng
Liu. 2021. {`{RoFormer}: Enhanced Transformer with Rotary Position
Embedding'}. \emph{arXiv Preprint arXiv:2104.09864}.

\bibitem[\citeproctext]{ref-var2024}
Tian, Keyu, Yi Jiang, Zehuan Yuan, Bingyue Peng, and Liwei Wang. 2024.
{`Visual Autoregressive Modeling: Scalable Image Generation via
Next-Scale Prediction'}. \emph{Advances in Neural Information Processing
Systems}.

\bibitem[\citeproctext]{ref-soundstream2021}
Zeghidour, Neil, Alejandro Luebs, Ahmed Omran, Jan Skoglund, and Marco
Tagliasacchi. 2021. {`SoundStream: An End-to-End Neural Audio Codec'}.
\emph{IEEE/ACM Transactions on Audio, Speech, and Language Processing}.

\end{CSLReferences}

\section*{Appendix A: Hyperparameters}\label{appendix-a-hyperparameters}
\addcontentsline{toc}{section}{Appendix A: Hyperparameters}

\begin{longtable}[]{@{}
  >{\raggedright\arraybackslash}p{(\linewidth - 14\tabcolsep) * \real{0.1250}}
  >{\raggedright\arraybackslash}p{(\linewidth - 14\tabcolsep) * \real{0.1250}}
  >{\raggedright\arraybackslash}p{(\linewidth - 14\tabcolsep) * \real{0.1250}}
  >{\raggedright\arraybackslash}p{(\linewidth - 14\tabcolsep) * \real{0.1250}}
  >{\raggedright\arraybackslash}p{(\linewidth - 14\tabcolsep) * \real{0.1250}}
  >{\raggedright\arraybackslash}p{(\linewidth - 14\tabcolsep) * \real{0.1250}}
  >{\raggedright\arraybackslash}p{(\linewidth - 14\tabcolsep) * \real{0.1250}}
  >{\raggedright\arraybackslash}p{(\linewidth - 14\tabcolsep) * \real{0.1250}}@{}}
\caption{Main settings per experiment. All runs: AdamW, weight decay
0.01, cosine schedule with warmup, batch 128, commitment
\(\beta = 0.25\), fine-level dropout 0.5.}\tabularnewline
\toprule\noalign{}
\begin{minipage}[b]{\linewidth}\raggedright
Run
\end{minipage} & \begin{minipage}[b]{\linewidth}\raggedright
Levels
\end{minipage} & \begin{minipage}[b]{\linewidth}\raggedright
Codes
\end{minipage} & \begin{minipage}[b]{\linewidth}\raggedright
Codec
\end{minipage} & \begin{minipage}[b]{\linewidth}\raggedright
Codec steps (lr)
\end{minipage} & \begin{minipage}[b]{\linewidth}\raggedright
Generator
\end{minipage} & \begin{minipage}[b]{\linewidth}\raggedright
Gen.~steps (lr)
\end{minipage} & \begin{minipage}[b]{\linewidth}\raggedright
Other
\end{minipage} \\
\midrule\noalign{}
\endfirsthead
\toprule\noalign{}
\begin{minipage}[b]{\linewidth}\raggedright
Run
\end{minipage} & \begin{minipage}[b]{\linewidth}\raggedright
Levels
\end{minipage} & \begin{minipage}[b]{\linewidth}\raggedright
Codes
\end{minipage} & \begin{minipage}[b]{\linewidth}\raggedright
Codec
\end{minipage} & \begin{minipage}[b]{\linewidth}\raggedright
Codec steps (lr)
\end{minipage} & \begin{minipage}[b]{\linewidth}\raggedright
Generator
\end{minipage} & \begin{minipage}[b]{\linewidth}\raggedright
Gen.~steps (lr)
\end{minipage} & \begin{minipage}[b]{\linewidth}\raggedright
Other
\end{minipage} \\
\midrule\noalign{}
\endhead
\bottomrule\noalign{}
\endlastfoot
E1 & 1,2,4,8,16,32 & 128 & 128 wide, 2+2 blocks & 3,000 (1e-3) & 6 × 256
& 3,000 (6e-4) & EMA 0.95; best step 1,500 \\
E4 & 1,3,9,27,54,81 & 32 & 2 encoder blocks & 3,000 (1e-3) & 8 × 256 &
5,000 (6e-4) & prefix 81 \\
E5 & 1,2,4,6,9,13,18, 26,37,50,67,81 & 16 & lookup encoder, 1 decoder
block, 256 wide & 3,000 (1e-3) & 8 × 256 & 4,000 (6e-4) & EMA 0.99 \\
E6 & as E5 & 16 & E5 + learned level 0, GroupNorm(8), dropout 0.3 &
3,000 (2e-4) & 8 × 256 & 4,000 (6e-4) & \\
E8 & as E5 & 16 & E6 codec & --- & 8 × 256 & 4,000 (6e-4) & corruption
0.30 → 0.05, compiled \\
E9 & as E5 & 16 & E6 codec & --- & 8 × 256 & 4,000 (6e-4) & as E8,
model-sampled \\
\end{longtable}

Later Sudoku runs reuse the frozen E6 codec and its 12-level, 16-code
hierarchy. All use an 81-cell prefix, eight attention heads, configured
generator seed 0 except E20 (seed 1), AdamW with weight decay 0.01, a
cosine learning-rate schedule and compiled generator training. Generator
dropout is 0.1 at width 256 and 0.2 at width 512; warmup is 200 steps at
width 256 except E19/E20, which use 1,000; width 512 also uses 1,000.
The codec is additional to the generator sizes below.

\begin{longtable}[]{@{}llllll@{}}
\caption{Later Sudoku generator settings. Steps are completed steps for
the reported results; corruption gives the first and last replacement
probabilities.}\tabularnewline
\toprule\noalign{}
Run & Generator & Steps & Batch & Learning rate & Corruption \\
\midrule\noalign{}
\endfirsthead
\toprule\noalign{}
Run & Generator & Steps & Batch & Learning rate & Corruption \\
\midrule\noalign{}
\endhead
\bottomrule\noalign{}
\endlastfoot
E10 & 8 × 512 & 24,000 & 64 & 3e-4 & 0.30 → 0.05 \\
E11 & 8 × 256 & 32,000 & 64 & 6e-4 & 0.30 → 0.05 \\
E12 & 8 × 512 & 50,776 & 64 & 3e-4 & 0.50 → 0.10 \\
E13 & 8 × 256 & 32,000 & 64 & 6e-4 & 0.50 → 0.10 \\
E14 & 8 × 256 & 128,000 & 64 & 6e-4 & 0.30 → 0.05 \\
E15 & 8 × 256 & 128,000 & 64 & 6e-4 & 0.50 → 0.10 \\
E16 & 8 × 256 & 32,000 & 64 & 6e-4 & 0.50 → 0.10 \\
E17 & 8 × 256 & 32,000 & 64 & 6e-4 & 0.50 → 0.10 \\
E18 & 8 × 256 & 128,000 & 64 & 6e-4 & 0.50 → 0.10 \\
E19 pilot & 8 × 256 & 23,895 & 128 & 6e-4 & 0.50 → 0.10 \\
E19 funded & 8 × 256 & 128,000 & 128 & 6e-4 & 0.50 → 0.10 \\
E20 & 8 × 256 & 128,000 & 128 & 6e-4 & 0.50 → 0.10 \\
\end{longtable}

The base data are 200,000 training puzzles and 1,000 validation puzzles,
with split seed 0 and 28--45 givens. E15 adds one million greedy-unique
puzzles for each seed 201--204; E16 adds one million for each seed
205--206. E17 adds the unique subset of one million random-removal
candidates for each seed 101--104; its final 2,000 steps are a
weights-only continuation with reset optimizer and random state (Section
6.5). E12's recorded run completed 50,776 steps; the repository's
current E12 configuration specifies 200,000, so running that
configuration unchanged does not reproduce the reported exposure. E18
adds no unique-puzzle corpus. Both E19 runs and E20 add one million
greedy-unique puzzles per seed 201--208; their separate schedules use
full-state snapshots every 1,000 updates. The original million-step E19
configuration remains unstarted and is not either reported result. Full
settings and checkpoint prerequisites are in the corresponding
\texttt{configs/sudoku-e*.json} files and \texttt{docs/EXPERIMENTS.md}.

\section*{Appendix B: Environment}\label{appendix-b-environment}
\addcontentsline{toc}{section}{Appendix B: Environment}

Apple M5 Pro, 64 GB unified memory, macOS 26.5; MLX 0.32.3, NumPy,
Python 3.12.14 managed with mise and uv. Figures: matplotlib. Paper:
Pandoc and Typst. E18: RTX 3080, Python 3.12.15; both E19 runs and E20:
NVIDIA A40, Python 3.12.3, MLX 0.32.3 CUDA. Preflight and contention
qualifications accompany each run in the experiment log.

\section*{Appendix C: Version history}\label{appendix-c-version-history}
\addcontentsline{toc}{section}{Appendix C: Version history}

{\def\LTcaptype{none} % do not increment counter
\begin{longtable}[]{@{}
  >{\raggedright\arraybackslash}p{(\linewidth - 4\tabcolsep) * \real{0.3333}}
  >{\raggedright\arraybackslash}p{(\linewidth - 4\tabcolsep) * \real{0.3333}}
  >{\raggedright\arraybackslash}p{(\linewidth - 4\tabcolsep) * \real{0.3333}}@{}}
\toprule\noalign{}
\begin{minipage}[b]{\linewidth}\raggedright
Version
\end{minipage} & \begin{minipage}[b]{\linewidth}\raggedright
Date
\end{minipage} & \begin{minipage}[b]{\linewidth}\raggedright
Changes
\end{minipage} \\
\midrule\noalign{}
\endhead
\bottomrule\noalign{}
\endlastfoot
0.1 & 2026-10-23 & First draft: method, implementation, TinyWorld and
Sudoku results E0--E6. \\
0.2 & 2026-10-23 & Plain-language introduction, glossary, formulas
explained with a running example, experiment journey (E0--E9),
hyperparameter appendix, NeurIPS-style layout, HTML, EPUB, LaTeX, JATS
and DOCX editions. \\
0.3 & 2026-10-23 & E8 results: conditioning corruption gives the first
solved boards. \\
0.4 & 2026-10-23 & E11: four times the training data gives 55\% of easy
boards; first run on NVIDIA. \\
0.5 & 2026-10-24 & E13 (stronger corruption) and E14 (16× data): 87\% of
easy boards. \\
0.6 & 2026-10-24 & E15 combines data, stronger corruption and unique
puzzles: 96\% of easy boards. Layout checks on every build. \\
0.7 & 2026-10-24 & Full-state pause/resume and recovery checks; E17's
weights-only continuation solves 71.25\% of easy boards, with the
optimizer-reset qualification. \\
0.8 & 2026-10-24 & Release audit: codec/system parameter accounting,
budget and sampling qualifications, prior-art scope and simplified
visual identity. \\
0.9 & 2026-10-24 & Submission preparation: public-data figure rebuild
command, qualified memory-bandwidth discussion, later Sudoku settings
and AI assistance disclosure. \\
0.10 & 2026-10-25 & Completed E18/E19 results, checkpoint provenance,
A40 timings, uncertainty and explicit budget differences. \\
0.11 & 2026-10-27 & Completed E20 seed replication; two-seed assessment,
hard-board failure accounting, explicit provenance limits and a gated
text-research direction. \\
\end{longtable}
}

\end{document}
